\subsection{\texorpdfstring{EXAMPLES OF \(\mathfrak{S}\)-CONVERGENCE}{EXAMPLES OF S-CONVERGENCE}}

I. Uniform convergence in a subset of X. Let A be a subset of X and take \(\mathfrak{S} = \{\mathrm{A}\}\) . The uniformity (resp. topology) of \(\mathfrak{S}\) -convergence is then called the uniformity (resp. topology) of uniform convergence in A; if a filter \(\Phi\) on \(\mathcal{F}_{\mathfrak{S}}(\mathrm{X};\mathrm{Y})\) converges to \(u_0\) , it is said to converge to \(u_0\) uniformly in A. When \(\mathrm{A} = \mathrm{X}\) we recover the structure of uniform convergence defined in no. 1.

\begin{enumerate}
\def\labelenumi{\Roman{enumi}.}
\setcounter{enumi}{1}
\tightlist
\item
  Pointwise convergence in a subset of X. Let A be a subset of X, and take S to be the set of all subsets of X which consist of a single point belonging to A (by Remark 2 of no. 2 it comes to the same thing if we take S to be the set of all finite subsets of A). The uniformity (resp. topology) of S-convergence is then called the uniformity (resp. topology) of pointwise convergence in A; if a filter Φ on \(\mathcal{F}_{\mathfrak{S}}(\mathbf{X};\mathbf{Y})\) converges to \(u_0\) , it is said to converge to \(u_0\) pointwise in A. This will be the case if and only if, for each \(x\in \mathbf{A}\) , \(u_0(x)\) is a limit of \(u(x)\) with respect to the filter Φ.
\end{enumerate}

In particular, when A = X, the uniformity (resp. topology) of pointwise convergence in X is called simply the uniformity (resp. topology) of pointwise convergence; the uniform space obtained by endowing \(\mathcal{F}(X; Y)\) with this structure is denoted by \(\mathcal{F}_{s}(X; Y)\) . Note that the topology of pointwise convergence is just the product topology on \(Y^{X}\) and therefore depends only on the topology of Y, and not on its uniform structure.

\begin{enumerate}
\def\labelenumi{\Roman{enumi}.}
\setcounter{enumi}{2}
\tightlist
\item
  Compact convergence. Suppose that X is a topological space, and take S to be the set of all compact subsets of X. The uniformity (resp. the topology) of S-convergence is then called the uniformity (resp. the topology) of compact convergence, and the uniform space obtained by endowing \(\mathcal{F}(X;Y)\) with this uniformity is denoted by \(\mathcal{F}_{c}(X;Y)\) . The structure of compact convergence is coarser than that of uniform convergence, and the two coincide if X is compact; also it is finer than the structure of pointwise convergence, and these two coincide if X is discrete.
\end{enumerate}

If X is a uniform space we can define on \(\mathcal{F}(X;Y)\) the uniformity of precompact convergence by taking S to be the set of all precompact subsets of X. Again, if X is a metric space we may take S to be the set of all bounded subsets of X; the uniformity of S-convergence is then called the uniformity of bounded convergence.

